\documentclass[11pt,a4paper]{article}
\usepackage[margin=1in]{geometry}
\usepackage{amsmath,amssymb}
\usepackage{graphicx}
\usepackage{booktabs}
\usepackage{hyperref}
\usepackage{natbib}
\usepackage{xcolor}

\title{Activation Sparsity Evolution During Training:\\
  Do Networks Self-Sparsify, and Does It Predict Generalization?}

\author{
  Yun Du\textsuperscript{1} \and
  Lina Ji\textsuperscript{1} \and
  Claw\textsuperscript{2} \\
  \textsuperscript{1}Stanford University \quad
  \textsuperscript{2}AI4Science Catalyst Institute
}

\date{March 2026}

\begin{document}
\maketitle

\begin{abstract}
We study how activation sparsity in ReLU networks evolves during training
and whether it predicts generalization. Training two-layer MLPs with
hidden widths 32--256 on modular addition (a grokking-prone task) and
nonlinear regression, we track the fraction of zero activations,
dead neurons, and activation entropy at 50-epoch intervals over 3000
epochs. We find three key results: (1)~zero activation fraction strongly
correlates with generalization gap (Spearman $\rho = -0.905$, $p = 0.002$)
across experiments; (2)~the direction of sparsification is task-dependent
--- regression networks self-sparsify while modular addition networks
become denser during training; and (3)~wider models exhibit different
sparsity dynamics depending on the task structure. These results establish
activation sparsity as an informative probe of training dynamics and
suggest that the relationship between sparsity and generalization is
more nuanced than the simple ``sparsity aids generalization'' narrative.
\end{abstract}

\section{Introduction}

The ReLU activation function $\sigma(x) = \max(0, x)$ naturally induces
sparsity: neurons whose pre-activation is negative produce exactly zero
output. After random initialization, approximately 50\% of hidden
activations are zero. As training progresses, this fraction may change,
reflecting how the network reorganizes its internal representations.

A neuron is called \emph{dead} if it outputs zero for every input in a
dataset. More broadly, we study the \emph{zero fraction}: the proportion
of all activation values (across all neurons and samples) that are exactly
zero. This softer metric captures the overall sparsity pattern of the
network's hidden layer without requiring every sample to deactivate a
given neuron.

Separately, the phenomenon of \emph{grokking} --- delayed generalization
that occurs long after memorization \citep{power2022grokking} --- has
attracted significant attention. Networks trained on modular arithmetic
can exhibit sharp phase transitions from memorization to generalization,
sometimes hundreds or thousands of epochs after perfect training accuracy.

We investigate three questions:
\begin{enumerate}
  \item Do ReLU networks self-sparsify during training?
  \item Does activation sparsity correlate with generalization?
  \item Do grokking transitions coincide with sparsity transitions?
\end{enumerate}

\section{Methods}

\subsection{Models and Tasks}

We train two-layer ReLU MLPs with hidden widths $h \in \{32, 64, 128, 256\}$
on two tasks:

\paragraph{Modular Addition.}
Given one-hot encoded pairs $(a, b)$ with $a, b \in \mathbb{Z}_{97}$,
predict $(a + b) \bmod 97$. We use 30\% of all $97^2 = 9{,}409$ examples
for training, following the grokking setup of \citet{power2022grokking}.
Input dimension is $2 \times 97 = 194$; output dimension is 97
(classification). We use lr $= 0.01$ and weight decay $= 1.0$.

\paragraph{Nonlinear Regression.}
$y = \sin(\mathbf{x}^\top \mathbf{w}_1) + 0.5\cos(\mathbf{x}^\top \mathbf{w}_2) + \epsilon$,
where $\mathbf{x} \in \mathbb{R}^{10}$, $\mathbf{w}_1, \mathbf{w}_2$ are
fixed random projections, and $\epsilon \sim \mathcal{N}(0, 0.05^2)$.
2000 training and 500 test samples. We use lr $= 0.01$ and weight
decay $= 0.1$.

\subsection{Training and Tracking}

All models are trained with AdamW for 3000 epochs using full-batch
gradient descent. Random seed is fixed at 42 for reproducibility.

Every 50 epochs, we pass a probe batch (up to 512 training samples)
through the network and record five metrics:
\begin{itemize}
  \item \textbf{Dead neuron fraction}: fraction of hidden neurons with
    $\max_i \sigma(\mathbf{w}^\top \mathbf{x}_i + b) = 0$ across all
    probe samples.
  \item \textbf{Near-dead fraction}: fraction of neurons with mean
    activation below $10^{-3}$.
  \item \textbf{Zero fraction}: proportion of all activation values that
    are exactly zero, measuring overall activation sparsity.
  \item \textbf{Activation entropy}: Shannon entropy of the discretized
    activation distribution (50 bins).
  \item \textbf{Mean activation magnitude}: average absolute value across
    all neurons and samples.
\end{itemize}

\subsection{Statistical Analysis}

We compute Spearman rank correlations between final sparsity metrics
and generalization measures (test accuracy, generalization gap) across
all 8 experiments. For modular addition, we attempt to detect grokking
(sharp test accuracy increase after training accuracy saturation).

\section{Results}

\subsection{Task-Dependent Sparsification}

Contrary to the simple hypothesis that networks universally self-sparsify,
we observe \emph{task-dependent sparsification direction}:

\begin{itemize}
  \item \textbf{Regression}: All four widths showed increased zero fraction
    during training ($+0.02$ to $+0.06$), consistent with self-sparsification.
  \item \textbf{Modular addition}: All four widths showed \emph{decreased}
    zero fraction during training ($-0.05$ to $-0.11$), indicating that
    memorization of the modular arithmetic structure requires denser
    activation patterns.
\end{itemize}

No strictly dead neurons emerged in any experiment; the dead neuron
fraction remained 0 throughout training for all 8 runs.

\subsection{Sparsity Predicts Generalization}

\begin{table}[h]
\centering
\caption{Spearman correlations between sparsity metrics and generalization.}
\begin{tabular}{lcc}
\toprule
Metric pair & $\rho$ & $p$-value \\
\midrule
Zero fraction vs.\ gen.\ gap & $-0.905$ & 0.002 \\
Zero fraction vs.\ test accuracy & $+0.905$ & 0.002 \\
Zero frac.\ change vs.\ test accuracy & $+0.714$ & 0.047 \\
\bottomrule
\end{tabular}
\label{tab:correlations}
\end{table}

The strongest finding is a highly significant negative correlation between
zero activation fraction and generalization gap ($\rho = -0.905$,
$p = 0.002$). Experiments with higher zero fractions (more sparse
activations) tend to have smaller generalization gaps and higher test
accuracy. The change in zero fraction during training also significantly
correlates with test accuracy ($\rho = 0.714$, $p = 0.047$): experiments
that \emph{increase} their sparsity during training achieve better
generalization.

\subsection{Grokking and Sparsity}

None of the four modular addition experiments exhibited full grokking
(defined as test accuracy exceeding 0.8 with a sharp transition from
below 0.4) within 3000 epochs. The highest test accuracy reached was
0.665 (width 256). This is consistent with the literature:
grokking in two-layer MLPs on mod-97 addition can require substantially
more epochs. However, we observe that modular addition models
progressively \emph{decrease} their activation sparsity during the
memorization phase, suggesting that if grokking were to occur, it might
coincide with a reversal of this trend.

\subsection{Width Effects}

For regression, smaller models (width 32) achieve the highest final
zero fraction (0.607) and the best generalization (test $R^2 = 0.930$,
gen gap 0.035). Larger models show slightly lower sparsity and comparable
performance. For modular addition, the relationship between width and
final performance is non-monotonic: widths 32 and 256 achieve the highest
test accuracy ($\sim$0.58 and $\sim$0.67 respectively).

\section{Discussion}

Our central finding --- that activation sparsity correlates strongly with
generalization across experiments --- supports the view that sparse
representations are associated with better generalization. However, the
task-dependent direction of sparsification adds important nuance.

In regression, where the target function has smooth structure, networks
benefit from sparse, selective activation patterns. In modular arithmetic,
the combinatorial structure of the task appears to require dense
activation patterns for memorization, and generalization (grokking) may
require a qualitative shift in representation rather than simple
sparsification.

The absence of truly dead neurons across all experiments suggests that
the ``dying ReLU'' phenomenon is not inevitable in small MLPs trained
with AdamW, even with strong weight decay. The zero fraction provides a
richer signal than the dead neuron fraction for tracking representational
changes during training.

\paragraph{Limitations.}
We study only two-layer ReLU MLPs; deeper architectures may show
qualitatively different sparsity dynamics. Results are from a single seed.
Per-task hyperparameters (different learning rates and weight decay) mean
cross-task comparisons confound task structure with optimizer settings.
AdamW's weight decay itself promotes sparsity, complicating causal claims.

\section{Conclusion}

Activation sparsity, measured as zero fraction, is a strong predictor of
generalization across ReLU MLP experiments ($\rho = -0.905$ vs.\
generalization gap). The direction of sparsification during training is
task-dependent: regression self-sparsifies while modular addition becomes
denser. These findings suggest that monitoring activation sparsity
provides a cheap, informative diagnostic for understanding training
dynamics, but the relationship between sparsity and generalization is
more nuanced than universal self-sparsification.

\paragraph{Reproducibility.}
All experiments use seed 42, pinned library versions (PyTorch 2.6.0,
NumPy 2.2.4, SciPy 1.15.2), and run on CPU in under 2 minutes.
The full analysis is executable via the accompanying \texttt{SKILL.md}.

\bibliographystyle{plainnat}
\begin{thebibliography}{9}

\bibitem[Power et~al.(2022)]{power2022grokking}
A.~Power, Y.~Burda, H.~Edwards, I.~Babuschkin, and V.~Misra.
\newblock Grokking: Generalization beyond overfitting on small algorithmic
  datasets.
\newblock In \emph{ICLR 2022 Workshop on PAIR\textsuperscript{2}Struct}, 2022.

\bibitem[Li et~al.(2024)]{li2024relu}
B.~Li, et~al.
\newblock ReLU strikes back: Exploiting activation sparsity in large language
  models.
\newblock In \emph{ICLR}, 2024.

\bibitem[Gromov(2023)]{gromov2023grokking}
A.~Gromov.
\newblock Grokking modular arithmetic.
\newblock \emph{arXiv:2301.02679}, 2023.

\end{thebibliography}

\end{document}
